\section{行列式与体积：线性变换是放大、翻转，还是压扁空间}
\label{sec:03-Linear-Algebra/04-Determinants/01}

拿一个具体的矩阵：

\[
A=\begin{pmatrix}a&c\\b&d\end{pmatrix},
\]

它的两列为

\[
a_1=\begin{pmatrix}a\\b\end{pmatrix},\qquad
a_2=\begin{pmatrix}c\\d\end{pmatrix}.
\]

单位正方形的标准基 \(e_1,e_2\) 被 \(A\) 分别映到 \(a_1,a_2\)，整个正方形变成由这两个列向量张成的平行四边形。

先不管数学定义，只做一件朴素的事：算出这个平行四边形的面积。

把平行四边形框进一个以 \((0,0)\) 和 \((a+c,\,b+d)\) 为对角的大矩形里。矩形的面积是 \((a+c)(b+d)\)。从矩形中切掉四个直角三角形的面积：

\[
\frac12ab+\frac12ab+\frac12cd+\frac12cd = ab+cd,
\]

再切掉两个小矩形（合计 \(2bc\)）。剩下

\[
(a+c)(b+d)-(ab+cd)-2bc = ad-bc.
\]

所以平行四边形的面积就是 \(|ad-bc|\)。如果交换两列的顺序，同样的计算只差一个符号。

这件事已经很接近正确答案了，但缺了一样东西：纯面积永远是正数或零，而有些变换会翻转平面——比如把右手坐标系变成左手坐标系。面积本身区分不了“有没有翻面”。

于是自然地引入带符号的版本：

\[
\det A = \det\begin{pmatrix}a&c\\b&d\end{pmatrix} = ad-bc.
\]

它的绝对值是平行四边形的面积；正负号记录了方向是否翻转。如果两列落在同一条直线上，平行四边形被压成线段，面积为零，公式也恰好给出零。

这就是二维情形的全部几何内容。下面把同样的直觉推到任意维度。

\subsection{三个要求唯一确定了行列式}

我们不从排列求和公式出发。那个公式只是计算结果，不是理解的起点。更好的做法是：先列出任何合理的“有向体积”必须满足的条件，看看这些条件本身是不是已经足够把公式锁死。

设 \(\det\) 是一个给方阵的列赋一个数的运算。让它符合我们对“体积”的直觉，至少需要三条：

\textbf{对每一列的线性性（多线性）。} 固定其他列不变，单独拉伸某一列——比如把 \(a_1\) 换成 \(\alpha a_1+\beta a_1'\)——体积也应该按同样的线性组合变化：

\[
\det(\alpha a_1+\beta a_1',a_2,\ldots,a_n)
= \alpha\det(a_1,a_2,\ldots,a_n)
+ \beta\det(a_1',a_2,\ldots,a_n).
\]

这很合理：把一条边拉长两倍或拆成两段的组合，体积也应该对应变化。

\textbf{两列相同则体积为零（交替性）。} 两条边重合，平行四边形退化成一线段，\(n\) 维体积消失。结合多线性，可从这一条推出交换任意两列会变号：

\[
\begin{aligned}
0 &= \det(a_1+a_2,\,a_1+a_2) \\
  &= \det(a_1,a_1) + \det(a_1,a_2) + \det(a_2,a_1) + \det(a_2,a_2) \\
  &= \det(a_1,a_2) + \det(a_2,a_1),
\end{aligned}
\]

其中 \(\det(a_1,a_1)=\det(a_2,a_2)=0\) 来自“两列相同”。于是

\[
\det(a_2,a_1) = -\det(a_1,a_2).
\]

交换两列等价于翻转一次方向，符号自然改变。

\textbf{单位矩阵对应体积 1（规范性）。} 单位矩阵不改变标准基，单位立方体映射后仍然是它自己，体积应该为 \(1\)：

\[
\det I = 1.
\]

这三条合在一起，在任意维度唯一地确定了行列式。唯一性本身是一个定理——需要借助排列的论证来证明——但作为起点，知道“只需要这三条就够了”比记住排列公式更重要。排列求和公式会在下一节的开头给出；它只是在满足三条要求之后的一种显式计算手段。

\subsection{高维中的含义}

对 \(A\in\R^{n\times n}\)，单位立方体

\[
\{x:0\leq x_i\leq 1\}
\]

被映成由 \(A\) 的列张成的 \(n\) 维平行多面体。它的 \(n\) 维有向体积就是 \(\det A\)；实际体积是 \(|\det A|\)。

更一般地，任何区域经过线性变换 \(x\mapsto Ax\) 后，体积都乘以 \(|\det A|\)。论证思路是把区域切成足够小的立方体，每个小立方体的体积都放大同一个因子 \(|\det A|\)，求和取极限，整个区域也放大同样的因子。这个“所有体积共享同一个缩放因子”的事实，在下一节推导 \(\det(AB)=\det A\,\det B\) 时是核心杠杆。

几种典型情形：

\begin{itemize}
\tightlist
\item
  \(|\det A|>1\)：总体积放大；
\item
  \(0<|\det A|<1\)：总体积缩小；
\item
  \(\det A<0\)：除缩放外还翻转了方向；
\item
  \(\det A=0\)：至少一个维度被压扁，\(n\) 维体积消失。
\end{itemize}

用一个例子把抽象变成数字：

\[
A=\begin{pmatrix}2&1\\0&3\end{pmatrix},
\qquad
\det A = 2\cdot3 - 0\cdot1 = 6.
\]

单位正方形变成了面积为 \(6\) 的平行四边形。上三角位置的 \(1\) 造成了剪切——平行四边形被推斜了——但面积完全没变。对角线上的 \(2\) 和 \(3\) 各自沿一个独立方向缩放，体积因子相乘得 \(6\)。行列式把剪切和缩放的效应干净地分开了：只有对角线方向的缩放才贡献到乘积里。

三维中的一个上三角矩阵同样爽快：

\[
A=\begin{pmatrix}2&1&1\\0&3&2\\0&0&4\end{pmatrix},
\qquad
\det A = 2\cdot3\cdot4 = 24.
\]

对角线之外的 \(1,1,2\) 都是剪切——在三个维度上各自推斜——但它们对体积的贡献为零。单位立方体经过这个矩阵后变成体积为 \(24\) 的平行六面体，每个对角元恰好对应一个主轴方向的缩放。三角矩阵的行列式永远只是对角元之积，这个事实不是公式记忆的结果，而是从“剪切不改变体积”与“各方向缩放独立相乘”两条几何直觉的直接推论。

\subsection{行操作对应哪些变化}

三类初等行操作对行列式的影响，可以直接从几何读出：

\begin{enumerate}
\tightlist
\item
  交换两行，方向翻转，行列式乘 \(-1\)；
\item
  一行乘以 \(c\)，该方向拉伸 \(c\) 倍，行列式也乘 \(c\)；
\item
  一行加上另一行的倍数，只做剪切，行列式不变。
\end{enumerate}

列操作有完全相同的规则，因为 \(\det(A^{\trans})=\det A\)（转置不改变行列式；这依赖于 \(\det(AB)=\det A\,\det B\) 以及行操作与列操作之间的对偶关系，下一节会给出完整推导）。

这组规则说明 Gaussian 消元不仅能判断可逆性，还能跟踪体积因子。把矩阵通过消元化为上三角后，行列式就是对角元之积，再按过程中使用过的交换和缩放次数做符号和倍数补偿。

走一个具体的例子。取

\[
A=\begin{pmatrix}1&2&3\\2&5&8\\0&3&6\end{pmatrix}.
\]

消元步骤和对应的行列式追踪：

第一步，第二行减去两倍第一行（剪切，行列式不变）：

\[
\begin{pmatrix}1&2&3\\0&1&2\\0&3&6\end{pmatrix},
\]

第二步，第三行减去三倍第二行（又一次剪切，行列式不变）：

\[
\begin{pmatrix}1&2&3\\0&1&2\\0&0&0\end{pmatrix},
\]

对角元之一的第三个是 \(0\)，所以 \(\det A=0\)。回头看原矩阵：第三列 \(3,8,6\) 恰好是第一列的 \(3\) 倍加上第二列的 \(2\) 倍——列线性相关。消元忠实地把这个依赖关系追踪到了底的零对角元上。行列式不需要是一个神秘的数字——它就是消元过程留下的体积追踪记录。

\subsection{行列式为零就是空间被压扁}

若 \(A\) 的列线性相关，它们只能张成低于 \(n\) 维的子空间。对应的 \(n\) 维平行多面体在一个低维子空间里，\(n\) 维体积必然为零，因此

\[
\det A = 0.
\]

反过来，若 \(\det A=0\)，列向量所围成的 \(n\) 维体积消失了，说明它们不能提供 \(n\) 个独立方向，矩阵不满秩。

对方阵，以下陈述都是同一件事的不同说法：

\[
\det A\neq0
\;\Longleftrightarrow\;
A\text{ 可逆}
\;\Longleftrightarrow\;
\rank(A)=n
\;\Longleftrightarrow\;
\mathcal N(A)=\{0\}.
\]

行列式用一个数把“可逆矩阵定理”里的多条等价条件浓缩了。但它并没有告诉我们零空间指向哪个具体方向，也没有告诉我们矩阵离奇异还有多远——这两条需要特征值和奇异值来回答。

\subsection{符号记录了什么}

二维和三维中的有序基可以分为两种手性：右手系和左手系。交换任意两个基向量会切换手性类别，也会让行列式变号。

旋转矩阵

\[
R_\theta=
\begin{pmatrix}
\cos\theta&-\sin\theta\\
\sin\theta&\cos\theta
\end{pmatrix}
\]

满足 \(\det R_\theta = \cos^2\theta+\sin^2\theta = 1\)：旋转保持面积，也不翻转方向。镜像矩阵

\[
F=\begin{pmatrix}1&0\\0&-1\end{pmatrix}
\]

满足 \(\det F = -1\)：大小不变，但平面被翻了一面。

正交矩阵的行列式只能是 \(1\) 或 \(-1\)。从 \(Q^{\trans}Q=I\)，利用下一节的乘法性质 \(\det(AB)=\det A\,\det B\) 以及 \(\det I=1\)，得到 \((\det Q)^2=1\Rightarrow\det Q=\pm1\)。所以正交变换要么是纯旋转（\(+1\)），要么是带一次镜像的旋转（\(-1\)）——没有中间地带。

\subsection{一个需要警惕的误解}

行列式的绝对值很小，不代表矩阵“接近奇异”。

取 \(0.01I_{100}\)——一个 \(100\times100\) 的单位矩阵每行乘了 \(0.01\)。它的行列式是 \(0.01^{100}=10^{-200}\)，小得惊人。但这个矩阵在所有方向上的缩放完全一致，条件数是 \(1\)——它是你能想到的“距离奇异最远”的矩阵之一。

行列式把所有方向的缩放乘在一起。在极高维度或整体尺度很小时，乘积自然可以极小而每个独立方向仍然稳健。判断数值敏感性需要奇异值和条件数，行列式不承担这个职责。

延伸阅读：MIT 18.06：Properties of Determinants。
